\section{Related work}
\label{sec:related}

The literature was mapped against the five hypotheses \emph{before} any
model was fit to real data: the map (\texttt{docs/PRIOR-ART.md}, SHA-256
\texttt{3f36fb62\ldots}) was committed on 2026-08-26 and hash-stamped, and
the hypothesis-status clarifications it forced were registered the same day
as Addendum~3 (Section~\ref{sec:design-prereg}). The verdicts below are
therefore dated evidence of what was known when the design was frozen, not
a reading of the results.

\subsection{Stochastic mortality models and the Cairns--Blake--Dowd programme}

\citet{lee1992modeling} model the log central death rate as
$\log m_{x,t} = \alpha_x + \beta_x \kappa_t$, estimate the bilinear term by
singular value decomposition and extrapolate $\kappa_t$ by a random walk with
drift; the prediction intervals for $\kappa_t$, and hence for life
expectancy, come with the model. \citet{girosi2007understanding} show that
the method is a multivariate random walk with drift on the vector of log
rates with a fixed age profile, so that all of its forecast uncertainty is
$\kappa_t$ uncertainty scaled by $\beta_x$, and that forecast age profiles
become non-smooth; both points matter for reading a Lee--Carter interval.
\citet{brouhns2002poisson} replace the least-squares fit by a Poisson
likelihood with the exposure as offset, the error law appropriate to death
counts, and \citet{brouhns2005bootstrapping} add a semiparametric bootstrap
that carries parameter uncertainty into the projection. Coherent
multi-population variants follow \citet{li2005coherent}, and
\citet{li2017coherent} replace the factor structure by a sparse vector
autoregression on mortality improvements.

\citet{cairns2006two} introduce the two-factor CBD model for post-60
mortality with explicit parameter uncertainty; \citet{renshaw2006cohort} add
a cohort term to Lee--Carter. \citet{cairns2009quantitative} compare eight
such models (M1--M8) on England and Wales and United States data by
Bayesian information criterion, robustness to the fitting period and
biological reasonableness; \citet{cairns2011mortality} take six of them
forward as \emph{density} forecasters and judge their fan charts ex ante;
\citet{dowd2010evaluating} test the residuals of the same six for goodness
of fit; and \citet{dowd2010backtesting} set out a backtesting framework for
their multi-period-ahead density forecasts and apply it ex post to English
and Welsh males. That last paper is the direct ancestor of the present
audit and deserves a precise description. It backtests six models---M1,
M2B, M3B, M5, M6 and M7, each in a parameter-certain and a
parameter-uncertain variant---by locating realised mortality within the
prediction intervals of forecasts made from earlier origins, one population
and one age range at a time. Three of its findings frame ours. It reports
that ``forecast performance tends to improve with higher ages'', the
ex-ante alternative to \Hfour. It notes, in a footnote, that ``the
positions of the individual observations within the prediction intervals
are not independent'' across horizons, and then measures marginal
positions anyway; the joint-path coverage of \Hthree\ is the quantity that
footnote points at and does not compute. And it lists life expectancy,
survival rates and annuity prices as metrics on which ``in principle,
backtests could be conducted'' and confines itself to the mortality rate;
\Hfive\ is that extension. Later members of the programme include the
Bayesian two-population model of \citet{cairns2011bayesian}, the general
procedure for model construction of \citet{hunt2014general} and the CBDX
workhorse of \citet{dowd2020cbdx}. \citet{booth2008mortality} review the
wider field. The \pkg{StMoMo} package \citep{villegas2018stmomo} implements
the family in a common generalised-linear framework and is the numerical
oracle for the classical arms of this study.

Two features of this programme frame ours. It established that stochastic
mortality models are to be judged as density forecasters, not point
forecasters. And its evaluation, apart from \citet{dowd2010backtesting}, was
in-sample or ex ante, on one or two populations, before neural models
existed and before any break in the modern data.

\subsection{Neural extensions of Lee--Carter}

\citet{hainaut2018neural} is an early neural analyser of mortality surfaces.
\citet{richman2021neural} embed country and sex into a network that
generalises the Lee--Carter structure to all HMD populations at once (train
1950--1999, test 2000--2016) and evaluate on mean squared error.
\citet{nigri2019deep} replace the random walk on $\kappa_t$ by an LSTM;
\citet{petnehazi2019mortality} ask whether recurrent networks beat
Lee--Carter, by point error. \citet{perla2021time} show that a shallow
convolutional network can be read as a generalisation of Lee--Carter and
that it transfers from the HMD to the United States Mortality Database.
\citet{scognamiglio2022calibrating} fit Lee--Carter and Poisson Lee--Carter
jointly across populations with a network that replicates their structure.
\citet{marino2023neural} deepen Lee--Carter and add interval estimation
for the future $\kappa_t$ by a bootstrap ensemble of LSTMs;
\citet{miyata2022extending} replace the two-stage procedure by a
variational autoencoder fit in one stage and obtain intervals without MCMC.

\citet{schnurch2022point} is the one neural mortality paper whose title
promises interval forecasts, and it is the closest prior work to ours. On
54 HMD populations at ages 60--89, trained to 2006 and tested on
2007--2016, they report the prediction-interval coverage probability (PICP)
and mean width at a nominal 95\% for Lee--Carter fit on ten and on twenty
years, an autoregressive-cohort model, and recurrent, feed-forward and
convolutional networks: coverages of 74.0, 74.3, 77.2, 86.0, 97.0 and 98.0
per cent respectively, with the neural intervals the widest. Their
intervals come from the variation across independently seeded trainings of
the same network, and their pass rule is one-sided (``minimize MPIW under
the constraint that PICP is at or above'' the nominal level), so
over-coverage is a pass. Two features of that study define what ours adds.
First, the uncertainty mechanism is held fixed by design: the authors
``rule out several existing methods from the literature because they would
require substantial structural changes, such as the insertion of dropout
layers for Monte Carlo dropout'', because doing so ``might decrease the
forecasting performance, the optimization of which is still our main
goal''. The 74-versus-98 per cent gap is therefore confounded between
architecture and mechanism, which is exactly what the crossed design of
Section~\ref{sec:design} separates. Second, their classical arms see ten or
twenty training years while the networks see every year of all 54
populations, so part of the classical under-coverage is a training-window
artefact; the expanding-origin rule of Section~\ref{sec:design-regimes}
gives every family the same window. Their Table~2 also shows the
point-versus-proper-score inversion of \Hone\ (mean squared error ranks the
networks first; Poisson deviance ranks Lee--Carter first), which they
attribute to exposure weighting rather than to distributional sharpness.

Adjacent non-neural machine-learning comparators are the gradient-boosted
multi-population model of \citet{miao2025gradient}, the tree-based analysis
of \citet{deprez2017machine}, the pension-fund application of
\citet{demelo2025actuarial}, and the Markov-switching Bayesian VAR on an
age-partitioned Lee--Carter of \citet{fu2025applying}, which is a Bayesian
rather than a neural route to uncertainty under structural change.

The evaluation profile of this lineage is documented in
Section~\ref{sec:intro}: point accuracy throughout; intervals plotted where
they exist; coverage computed without a sharpness penalty where it is
computed at all; proper scores rarely. Its stable-regime coverage record is
nevertheless informative, and it already settles part of \Htwo: in periods
without a break, classical intervals under-cover and bootstrap-ensemble
neural intervals over-cover \citep{schnurch2022point,marino2023neural}. The
registered ``under-cover for every family'' clause was therefore falsified
in the stable regime before our first fit, and the stable regime of this
paper is a replication arm (Section~\ref{sec:design-status}). The four
neural families in our grid (Section~\ref{sec:design-families}) are drawn
from this lineage.

\subsection{Bayesian and Gaussian-process mortality models}

\citet{ludkovski2018gaussian} model mortality surfaces by Gaussian-process
regression; \citet{huynh2021multi} extend this to multi-output GPs over many
HMD populations and, alone in our library, evaluate by CRPS as the score is
meant to be used. \citet{barigou2023bayesian} average Bayesian mortality
models by leave-future-out validation with the expected log predictive
density and CRPS across four countries, and find that ``the `optimal' model
depends on the country and the scoring rule studied'';
\citet{goes2025bayesian} report, on an England and Wales war-shock window,
a top-one inversion between CRPS and the point and log-score rankings.
These two papers, with \citet{schnurch2022point}, are why \Hone\ is prior
art and is carried as a harness-validity check rather than a finding
(Section~\ref{sec:design-status}). \citet{liu2020bayesian} fit a Bayesian
Poisson log-normal model with regularised time structure. These are the
models whose predictive distributions are honest by construction, and their
evaluation practice is the one we adopt for every family; a multi-output GP
is one of the ten families audited.

\subsection{COVID-era mortality modelling}

The pandemic literature in actuarial science is about \emph{estimation with}
the pandemic years, not \emph{evaluation across} them.
\citet{cairns2020impact} frame COVID-19 deaths as accelerated deaths with
implications for higher-age mortality; \citet{robben2022assessing} quantify
what one pandemic data point does to the calibration and projections of a
stochastic multi-population model; \citet{goes2025bayesian} extend
Lee--Carter with a vanishing jump so that pandemic years neither dominate the
trend nor have to be deleted; \citet{robben2024catastrophe} add catastrophe
risk to a multi-population model; \citet{vanberkum2025estimating} extend the
Li--Lee model to measure the impact on granular data;
\citet{nalmpatian2023advanced} combine generalised additive models with
machine learning to estimate the trend deviation; and
\citet{robben2025granular} model epidemic and temperature shocks by regime
switching. \citet{barigou2023bayesian}, as noted, simulate a pandemic
perturbation because real post-shock data did not yet exist, and remark
that adapting their validation to ``a change of regime'' is ``beyond the
scope'' of their paper. \citet{goes2025bayesian} state the constraint
directly: because the pandemic ``predominantly impacts the most recent
years of data'', an out-of-sample test was ``impractical'' and their COVID
analysis ``focuses solely on evaluating the in-sample fit''. None of these
papers evaluates prediction intervals that were fit before 2020 against
what happened in 2020--2024; with twenty populations now final through
2024, the constraint that stopped them is gone. That is the gap this audit
fills, and the papers above are what a modeller would reach for afterwards
to repair a failure we document; repairing it is not our contribution.

\subsection{Proper scoring rules, calibration and sharpness}

A scoring rule is proper if the expected score is optimised by the true
distribution and strictly proper if uniquely so; the logarithmic score and
the continuous ranked probability score are the two canonical examples
\citep{gneiting2007strictly}, and the interval (Winkler) score is the proper
score for a central prediction interval \citep{winkler1972decision}.
\citet{gneiting2007probabilistic} set out the paradigm that the present audit
adopts wholesale---maximise sharpness subject to calibration---together
with the probability integral transform (PIT) as the diagnostic of
probabilistic calibration, following the prequential principle of
\citet{dawid1984present}. \citet{czado2009predictive} give the randomised
PIT for count data, which we need because deaths are counts and predictive
samples tie. The decomposition of a proper score into uncertainty,
resolution and miscalibration goes back to \citet{murphy1973new};
\citet{pohle2020murphy} generalises it to all forecast types and
\citet{dimitriadis2021stable} give a stable estimator via isotonic
recalibration. \citet{dawid2014theory} survey the theory;
\citet{yanagisawa2023proper} extend proper scoring to survival-time
distributions, which is the natural target when the object of interest is a
lifetime rather than a rate.

\subsection{Conformal prediction}

Conformal prediction \citep{vovk2005algorithmic} produces prediction sets
with finite-sample marginal coverage under exchangeability of the calibration
and test residuals; split conformal \citep{lei2018distribution} makes it
cheap for any black-box regressor, and conformalised quantile regression
\citep{romano2019conformalized} makes the sets adaptive. Its relevance to a
time series is the question of what happens when exchangeability fails.
\citet{tibshirani2019conformal} reweight under covariate shift;
\citet{barber2023conformal} bound the coverage gap of weighted conformal
sets by a total-variation distance between calibration and test
distributions---a bound that a structural break is designed to make large.
\citet{xu2021conformal} propose ensemble batch prediction intervals (EnbPI)
for dynamic time series without exchangeability; \citet{gibbs2021adaptive}
adapt the miscoverage level online under distribution shift, and
\citet{gibbs2024conformal} extend this to arbitrary shifts. For multi-step
forecasts, where the object of interest is the whole path, joint coverage
is solved prior art in machine learning: \citet{stankeviciute2021conformal}
obtain simultaneous multi-horizon bands from a recurrent network by a
Bonferroni allocation, \citet{sun2024copula} calibrate a joint band through
a copula of the per-horizon scores, \citet{messoudi2021copula} build
copula-based multi-target sets, and \citet{diquigiovanni2022conformal} give
functional bands from a scaled supremum-norm score. \citet{sun2025conformal}
address change points directly. \citet{angelopoulos2023gentle} is the
standard primer.

We found no application of conformal prediction to mortality forecasting:
in the 52-paper corpus of the prior-art map, ``conformal'' has no hit in any
mortality text, and neither do ``Monte Carlo dropout'' or ``deep
ensemble''. Three conformal mechanisms enter our grid as \emph{audited
arms}, on the same footing as the native and ensemble mechanisms; the
joint-path arm is positioned against \citet{stankeviciute2021conformal}
and \citet{diquigiovanni2022conformal}, whose construction it uses, and we
claim the first \emph{measurement} of the marginal-to-joint coverage gap in
mortality forecasting, not the method. We make no claim that any conformal
arm fixes anything, and we deliberately do not invoke the finite-sample
guarantee under a break, for exactly the reason \citet{barber2023conformal}
give.

\subsection{Uncertainty mechanisms for neural networks}

Deep ensembles \citep{lakshminarayanan2017simple} estimate predictive
uncertainty from the disagreement of independently initialised networks;
Monte Carlo dropout \citep{gal2016dropout} reads dropout at prediction time
as approximate Bayesian inference. Both capture variability that the
training data can reveal; neither is designed to anticipate a regime the
training data never contained. Whether that matters for coverage at a break
is one of the questions the crossed design answers.

\subsection{Forecast-comparison inference}

\citet{diebold1995comparing} test equal predictive accuracy on a
loss-differential series; \citet{hansen2011model} construct the model
confidence set, the set of models not significantly outperformed at a given
confidence level; \citet{shang2018model} apply the latter to age-specific
mortality (Japanese national and prefectural data, retirement ages), the one
prior use we found in this setting. Because our shift regime is a single
common shock observed in twenty correlated populations, the effective
number of independent observations is nearer one than twenty, and ordinary
Diebold--Mariano standard errors overstate the information available; we
cluster by population, impose the null, and use the wild cluster bootstrap
of \citet{cameron2008bootstrap} with the six-point weights of
\citet{webb2023reworking}, which \citet{mackinnon2017wild} show remain
reliable at the small number of clusters we have. Rolling-origin evaluation
follows \citet{tashman2000out}; pre-registration follows the practice
advocated by \citet{nosek2018preregistration}.

\subsection{What is and is not prior art}

For the record, and because it was written down before any result existed,
the prior-art map's verdict on each hypothesis is: \Hone\ is prior art
\citep{barigou2023bayesian,goes2025bayesian,schnurch2022point}; the
stable-regime half of \Htwo\ is prior art and contradicts the registered
universality, while its shift-regime half is untested anywhere; \Hthree\ is
novel in mortality as a measurement and not novel as a method; \Hfour's
registered direction is contradicted by \citet{dowd2010backtesting}; and
\Hfive's construction is prior art while an ex-post coverage audit of the
derived quantities is not found. The white space that remains, and that
this paper occupies, is the crossed family-by-mechanism design, the
conformal arms, the marginal-to-joint gap, the derived-quantity coverage
audit, and the out-of-sample pandemic test the field declined for lack of
data. Two items are still unread at the time of writing and are named so
that their absence is visible: the online supplement of
\citet{schnurch2022point} (Section~C.3, implied annuity present values),
which is the one live threat to \Hfive's residual novelty if it audits
annuity-interval coverage rather than presenting point values, and the full
text of \citet{stankeviciute2021conformal}, cited above from its
description in later work.
